\SourcePage{121}{124}
\section{Matrix elements of vectors}

Return once more to a closed system of particles\footnote{Everything derived in
this section applies equally to a particle in a centrally symmetric field (more
broadly, to any situation in which the system's total angular momentum is
conserved).}. Let $f$ stand for some scalar physical quantity of this system
and $\hat f$ for its operator. Any
\SourcePage{122}{125}
scalar is invariant under a rotation of the coordinate system. Hence the
scalar operator $\hat f$ is unchanged by a rotation; that is, it commutes with
the rotation operator. We know, however, that the operator of an infinitesimal
rotation coincides, apart from a constant factor, with the angular-momentum
operator. Therefore
\begin{equation}
  \{\hat f,\hat{\mathbf L}\}=0.
  \tag{29.1}
\end{equation}

It follows from the commutation of $\hat f$ with the angular-momentum operator
that, for transitions between states with definite $L$ and $M$, the matrix of
$f$ is diagonal in these indices. Moreover, specifying $M$ determines only
the orientation of the system relative to the coordinate axes, while the
value of a scalar quantity does not depend on this orientation at all.
Consequently, the matrix elements $\langle n'LM|f|nLM\rangle$ do not depend
on $M$ (the letter $n$ provisionally denotes the collection of all quantum
numbers, other than $L$ and $M$, that specify the state of the system). A
formal proof follows from the commutation of $\hat f$ and $\hat L_+$:
\begin{equation}
  \hat f\hat L_+-\hat L_+\hat f=0.
  \tag{29.2}
\end{equation}
Take the matrix element of this equality for the transition
$n,L,M\to n',L,M+1$. Since the matrix of $L_+$ has only elements for
transitions $n,L,M\to n,L,M+1$, we find
\begin{align*}
 &\langle n',L,M+1|f|n,L,M+1\rangle
 \langle n,L,M+1|L_+|n,L,M\rangle={}\\
 &\qquad=\langle n',L,M+1|L_+|n',L,M\rangle
 \langle n',L,M|f|n,L,M\rangle,
\end{align*}
and, since the matrix elements of $L_+$ do not depend on the index $n$,
\begin{equation}
 \langle n',L,M+1|f|n,L,M+1\rangle
 =\langle n',L,M|f|n,L,M\rangle,
 \tag{29.3}
\end{equation}
which shows that all the elements $\langle n',L,M|f|n,L,M\rangle$ with
different $M$ (and identical remaining indices) are equal.

Applied to the Hamiltonian itself, this result gives the already known
independence of the stationary-state energy from $M$, that is, the
$(2L+1)$-fold degeneracy of the energy levels.

Next, let $\mathbf A$ be some vector physical quantity characterizing the
closed system. Under a rotation of the coordinate system (in particular, an
infinitesimal rotation, that is,
\SourcePage{123}{126}
under the action of the angular-momentum operator), the components of the
vector transform into one another. The commutation of the operators
$\hat L_i$ with the operators $\hat A_i$ must therefore again yield components
of the same vector $\hat A_i$. Which components occur can be found by noting
that, in the special case where $\mathbf A$ is the radius vector of a
particle, the result must reduce to (26.4). We thus obtain the commutation
rules
\begin{equation}
  \{\hat L_i,\hat A_k\}=ie_{ikl}\hat A_l.
  \tag{29.4}
\end{equation}

From these relations one can derive a number of results about the structure
of the component matrices of the vector $\mathbf A$
(\emph{M.~Born, W.~Heisenberg, P.~Jordan}, 1926). First, they make it possible
to find the \emph{selection rules}, which determine the transitions for which
matrix elements may be nonzero. We shall not, however, give the corresponding
rather cumbersome calculations here, since it will be shown below
(\secref{107}) that these rules are in fact direct consequences of the general
transformation properties of vector quantities and can be obtained from them
essentially without calculation. Here we merely state the rules without
derivation.

Matrix elements of every component of a vector can be nonzero only for
transitions in which the angular momentum $L$ changes by no more than one:
\begin{equation}
  L\to L,L\pm1.
  \tag{29.5}
\end{equation}
There is also an additional selection rule forbidding transitions between
any two states with $L=0$. This rule is an evident consequence of the complete
spherical symmetry of states with zero angular momentum.

The selection rules for the angular-momentum projection $M$ differ among the
components of the vector. Specifically, matrix elements may be nonzero for
transitions with the following changes of $M$:
\begin{equation}
 \begin{aligned}
  M&\to M+1 &&\text{for } A_+=A_x+iA_y,\\
  M&\to M-1 &&\text{for } A_-=A_x-iA_y,\\
  M&\to M   &&\text{for } A_z.
 \end{aligned}
 \tag{29.6}
\end{equation}
It is also possible to determine in general form the dependence of the matrix
elements of a vector on $M$. We again give these important and frequently
used formulas without derivation, since they too are special cases of more
general relations, applying to arbitrary tensor quantities, which will be
derived in \secref{107}.
\SourcePage{124}{127}
The nonzero matrix elements of $A_z$ are given by
\begin{equation}
 \begin{aligned}
 \langle n'LM|A_z|nLM\rangle
 &=\frac{M}{\sqrt{L(L+1)(2L+1)}}\langle n'L\Vert A\Vert nL\rangle,\\
 \langle n'LM|A_z|n,L-1,M\rangle
 &=\sqrt{\frac{L^2-M^2}{L(2L-1)(2L+1)}}
   \langle n'L\Vert A\Vert n,L-1\rangle,\\
 \langle n',L-1,M|A_z|nLM\rangle
 &=\sqrt{\frac{L^2-M^2}{L(2L-1)(2L+1)}}
   \langle n',L-1\Vert A\Vert nL\rangle.
 \end{aligned}
 \tag{29.7}
\end{equation}
Here the symbol
\[
  \langle n'L'\Vert A\Vert nL\rangle
\]
denotes a so-called \emph{reduced matrix element}, a quantity independent of
the quantum number $M$\footnote{The appearance in (29.7) and (29.9) of
denominators depending on $L$ accords with the general notation introduced in
\secref{107}. The usefulness of these denominators is seen, in particular, in the
simple form taken by (29.12) for the matrix elements of the scalar product of
two vectors.

One should regard the notation for the reduced matrix element as an indivisible entity
(in contrast to what was said about the matrix element notation (11.17)).}.
They are related by
\begin{equation}
 \langle n'L'\Vert A\Vert nL\rangle
 =\langle nL\Vert A\Vert n'L'\rangle^*,
 \tag{29.8}
\end{equation}
which follows directly from the Hermiticity of $\hat A_z$.

The matrix elements of $A_-$ and $A_+$ can be expressed through the same
reduced elements. The nonzero matrix elements of $A_-$ are
\begin{equation}
 \begin{aligned}
 \langle n',L,M-1|A_-|nLM\rangle
 &=\sqrt{\frac{(L-M+1)(L+M)}{L(L+1)(2L+1)}}
   \langle n'L\Vert A\Vert nL\rangle,\\
 \langle n',L,M-1|A_-|n,L-1,M\rangle
 &=\sqrt{\frac{(L-M+1)(L-M)}{L(2L-1)(2L+1)}}
   \langle n'L\Vert A\Vert n,L-1\rangle,\\
 \langle n',L-1,M-1|A_-|nLM\rangle
 &=-\sqrt{\frac{(L+M-1)(L+M)}{L(2L-1)(2L+1)}}
   \langle n',L-1\Vert A\Vert nL\rangle.
 \end{aligned}
 \tag{29.9}
\end{equation}
The matrix elements of $A_+$ need no separate formulas, since the reality of
$A_x$ and $A_y$ gives
\begin{equation}
 \langle n'L'M'|A_+|nLM\rangle
 =\langle nLM|A_-|n'L'M'\rangle^*.
 \tag{29.10}
\end{equation}
\SourcePage{125}{128}
Let us write out the formula that relates the matrix elements of the scalar
$\mathbf{AB}$ to the reduced matrix elements of the two vectors
$\mathbf A$ and $\mathbf B$. The calculation is conveniently made
by writing the operator $\hat{\mathbf A}\hat{\mathbf B}$ as
\begin{equation}
 \hat{\mathbf A}\hat{\mathbf B}
 =\frac12(\hat A_+\hat B_-+\hat A_-\hat B_+)+\hat A_z\hat B_z.
 \tag{29.11}
\end{equation}
The matrix of $\mathbf{AB}$, like that of any scalar, is diagonal in $L$ and
$M$. Calculation using (29.7)--(29.9) gives
\begin{equation}
 \langle n'LM|\mathbf{AB}|nLM\rangle
 =\frac1{2L+1}\sum_{n'',L''}
 \langle n'L\Vert A\Vert n''L''\rangle
 \langle n''L''\Vert B\Vert nL\rangle,
 \tag{29.12}
\end{equation}
where $L''$ takes the values $L,L\pm1$.

For reference, we give the reduced matrix elements of the vector
$\mathbf L$ itself. Comparing (29.9) and (27.12), we find
\begin{equation}
 \begin{gathered}
  \langle L\Vert L\Vert L\rangle=\sqrt{L(L+1)(2L+1)},\\
  \langle L-1\Vert L\Vert L\rangle
  =\langle L\Vert L\Vert L-1\rangle=0.
 \end{gathered}
 \tag{29.13}
\end{equation}

A quantity frequently encountered in applications is the unit vector
$\mathbf n$ along the radius vector of a particle. Let us find its reduced
matrix elements. It is sufficient to calculate, for example, the matrix
elements of $n_z=\cos\theta$ at zero angular-momentum projection, $m=0$. We
have
\[
 \langle l-1,0|n_z|l0\rangle
 =\int_0^\pi\Theta_{l-1,0}^*\cos\theta\cdot
 \Theta_{l0}\sin\theta\,d\theta
\]
with the functions $\Theta_{l0}$ from (28.11). Evaluation of the integral
gives\footnote{The calculation is performed by integrating by parts
$(l-1)$ times with respect to $d\cos\theta$. For the general formula for
integrals of this kind, see (107.14).}
\[
 \langle l-1,0|n_z|l0\rangle
 =\frac{il}{\sqrt{(2l-1)(2l+1)}}.
\]
The matrix elements for transitions $l\to l$ are zero (as they are for any
polar vector belonging to a single particle; see below, (30.8)). Comparison
with (29.7) now gives
\begin{equation}
 \langle l-1\Vert n\Vert l\rangle
 =-\langle l\Vert n\Vert l-1\rangle=i\sqrt l,
 \qquad \langle l\Vert n\Vert l\rangle=0.
 \tag{29.14}
\end{equation}
\SourcePage{126}{129}
\ProblemHeading

Average the tensor $n_in_k-(1/3)\delta_{ik}$ (where $\mathbf n$ is the unit
vector along the particle's radius vector) over a state with a definite value
of $|\mathbf l|$ but no preferred orientation (that is, $l_z$ is
unspecified).

\SolutionHeading The required mean value is an operator that can be expressed
only in terms of $\hat{\mathbf l}$. We seek it in the form
\[
 \overline{n_in_k}-\frac13\delta_{ik}
 =a\left[\hat l_i\hat l_k+\hat l_k\hat l_i
 -\frac23\delta_{ik}l(l+1)\right];
\]
this is the most general form of a symmetric second-rank tensor with zero
trace constructed from the components of $\hat{\mathbf l}$. To determine the
constant $a$, multiply this equality on the left by $\hat l_i$ and on the
right by $\hat l_k$, with summation over $i$ and $k$. Since the vector
$\mathbf n$ is perpendicular to the vector
$\hbar\hat{\mathbf l}=\hat{\mathbf r}\times\hat{\mathbf p}$, we have
$n_i\hat l_i=0$. We replace the product
$\hat l_i\hat l_i\hat l_k\hat l_k=(\hat{\mathbf l}^2)^2$ by its eigenvalue
$l^2(l+1)^2$, and transform the product
$\hat l_i\hat l_k\hat l_i\hat l_k$ using the commutation relations (26.7) as
follows:
\begin{align*}
 \hat l_i\hat l_k\hat l_i\hat l_k
 &=\hat l_i\hat l_i\hat l_k\hat l_k
 -ie_{ikl}\hat l_i\hat l_l\hat l_k\\
 &=(\hat{\mathbf l}^2)^2
 -\frac i2e_{ikl}\hat l_i(\hat l_l\hat l_k-\hat l_k\hat l_l)\\
 &=(\hat{\mathbf l}^2)^2
 +\frac12e_{ikl}e_{lkm}\hat l_i\hat l_m
 =(\hat{\mathbf l}^2)^2-\hat{\mathbf l}^2\\
 &=l^2(l+1)^2-l(l+1)
\end{align*}
(we have used $e_{ikl}e_{mkl}=2\delta_{im}$). Straightforward simplification
then gives
\[
  a=-\frac1{(2l-1)(2l+3)}.
\]
